\subsection{A broader bounded quotient screen}
The new search permits both rows of a retained control to be synthesized jointly,
including intermediate returns to that retained control.  Earlier independent-row
screens stopped when such a return occurred.  A partial assignment of one row
is rejected only if a required local trajectory fails before consulting the
unassigned row.  Surviving pairs of assignments are then checked together.
This is a discovery procedure; no smaller machine survived it.

The sufficient mathematical criterion is the following.  Let $M$ be the parent
machine, $M'$ its proposed quotient, and $E$ the specified configuration encoding.
For each working control $q$, let $A_q$ be a finite set of radius-$r$ windows
that overapproximates all parent configurations reachable from the initial tape.
Suppose that for every window in every $A_q$, the parent and proposed quotient
have positive finite trajectory segments, of lengths between $1$ and $B$, whose
endpoints satisfy $E$.  Both segments must be checked on the complete input
window, and must preserve every exterior cell.  The encoded initial configuration
must be the quotient's initial configuration.  Halt is terminal and is related
only to halt.  Then the two machines have the same blank-tape halting behavior.

Indeed, concatenate the certified segments along the parent computation.  At each
join the parent again lies in the overapproximated reachable set, so the next
certificate applies.  Positive progress on both sides prevents infinitely many
parent steps from being represented by finitely many quotient steps.  Conversely,
a terminal halt cannot be crossed by a certified segment, so a finite halting
computation on either side must end at a halting join.  A parent halting trace of
length $t$ gives a quotient halting trace of length at most $Bt$; the converse has
the same bound.  The construction and these finite inductions are formalizable
in PA.  Here this general statement describes the search criterion, rather than
claiming a newly kernel-checked theorem or a successful quotient.

The identity encoding identifies one removed control with one retained control.
Additional experiments used a fixed two-cell head displacement or a toggle of
the cell just written on entry, whenever all external arrivals to the removed
control had the same direction.  Incoming transition actions were changed
accordingly, and the endpoint relation included the displacement and toggle.
The two-sided screen may advance the parent as well as the quotient before
checking the join.  Its bounds remain genuine restrictions: it excludes valid
simulations that require larger windows, longer segments, more changed controls,
or a different tape encoding.
